\section{秦：统一来得很快，帝国来得更快}

战国末年，六国还在，地图上没有哪一国自愿承认自己将被抹去。可秦已经拥有一种不同的节奏：关中能提供相对稳定的腹地，县邑和军功把人口、土地与兵役接得更紧，长期战争又使它比对手更熟悉怎样把一次胜利变成下一次战争的准备。

\eventanchor{QIN-0230-UNIFICATION}{秦·统一六国}

前230年，韩先亡；赵、魏、楚、燕、齐随后相继失去独立地位。每一国的最后一战都不同：魏都大梁面临水攻，楚的抵抗更艰难，齐长期孤立后在前221年降秦。把这段历史压成“秦横扫六国”，会掩盖统一真正的形状：它是十余年间一连串让对手无法同时互救的战争、外交、道路、补给和内部判断。

统一之后，秦没有停在战场上。度量衡、文字、道路、郡县、法律和巡行被迅速推向新得到的土地。里耶秦简中的县级文书，让后人看见这一速度并非传说：原楚地已经被写进秦的行政日常。帝国的厉害之处，不只在它能打到哪里，而在它能多快让远方开始使用同一套文书、名目和命令。

\begin{sourcenote}[统一不是单一人物的一步棋]
\sourcebook{史记·秦始皇本纪}和《六国年表》保存了六国覆灭的基本次序；各国世家补足不同角度。王翦“六十万”、王贲灌大梁等细节可见于史书，但数字和人物对白仍需审慎。里耶秦简、巡游刻石和度量衡器物能够旁证统一后行政与标准化的实际推进。
\end{sourcenote}

\subsection{帝国的日常，先于帝国的崩溃}

秦统一后最容易被记住的是严刑、徭役和焚书；但如果只写这些，便无法解释它为何能在短时间内统治广阔地域。秦制真正改变的是日常：户籍怎样被记录，仓储怎样被核验，文书怎样来往，地方官怎样报告，度量衡怎样让不同地区的人进入同一套交换与征发体系。

这也是它最危险的地方。战国时代，国家把社会资源推到战争前线；统一以后，同一套动员机器仍以极高强度运转。道路、工程、巡行、军役和法律并不是孤立的暴政清单，而是一个中心不愿降低速度时，落在无数家庭身上的连续要求。

\eventanchor{QIN-0213-EMPIRE}{秦·高强度治理}

前213年的禁书政策和前212年“坑儒”故事，后来成为秦的黑色象征。汉代史家写它们时，已经站在一个以崇儒自居的新王朝里。可以确认的是，秦确有禁书和严厉处置的政策；不能轻易确认的是执行范围、被害者身份和后世叙事中被不断放大的数字。把“坑儒”直接写成系统清除现代意义上的儒家，也会遮住当时“儒”“方士”“诸生”等称谓本身的复杂性。

\subsection{大泽以后}

\eventanchor{QIN-0209-FALL}{秦·起事与覆亡}

前209年，大泽起事成为秦末裂缝露出的时刻。它不是所有不满的起点，而是许多已经在路上的人突然发现中心可以被挑战的信号。陈胜吴广、项羽、刘邦和各地旧贵族、地方武装先后进入同一场崩解；秦的文书和道路仍在运行，秦的命令却越来越难以抵达能执行的人。

到前206年，秦亡。它只存在了很短的时间，却留下了很长的制度影子。汉不会回到周的封国世界，也不会简单废除郡县、法律和统一标准。真正的问题变成：怎样继承秦所创造的国家能力，又不让国家能力以秦末那样的强度压垮自己。

\begin{textwindow}[贾谊《过秦论》]
“仁义不施而攻守之势异也。”

贾谊把秦的兴亡压缩成这一句。它有巨大的解释力量，也有汉初政治论说的立场：贾谊关心的是汉该怎样避免成为第二个秦。把它放在秦亡之后，不是要让它替代所有原因，而是让读者看见，秦亡不久，后人已经开始用秦的命运讨论新帝国该怎样统治。
\end{textwindow}

\begin{turningpoint}[制度得失：秦留下的不是一个答案]
秦用战国形成的动员能力完成统一，解决了长期分裂、标准不一和诸侯战争的问题；代价是中心把统一后的社会仍按高强度战时逻辑对待。秦亡证明这种逻辑难以长期承受，却没有证明统一帝国本身不能存在。汉以后所有王朝都要在秦留下的能力与代价之间重新作答。
\end{turningpoint}
